% ---------------------------------------------------------------------------
% Pillar 2: Zero-Variance Fraction diagnostic
% Iter 26 — cross-library, cross-scale elevation vs AERO (arXiv 2602.14338).
% Materialised from scripts/zvf_diagnostic.py outputs:
%   experiments/results/zvf_summary.tsv
%   experiments/results/zvf_by_library.tsv
%   experiments/results/zvf_failure_correlation.tsv
%   figures/zvf_vs_failure.pdf
%   figures/zvf_by_library.pdf
% ---------------------------------------------------------------------------

\subsection{Cross-Library ZVF Diagnostic vs AERO}
\label{sec:zvf-cross-library}
We elevate ZVF from a per-experiment sanity check to a
\emph{cross-library} diagnostic by aggregating the per-step ZVF emitted
by every mitigation method we ran on a single managed runtime (Tinker
rollout workers, identical reward parser, identical $\beta{=}0$ KL, and
identical chat template). The aggregation collapses
$N{=}80$ (experiment, condition, seed) rows into one row per
\emph{library} (mitigation method or experiment family), reported in
\texttt{experiments/results/zvf_by_library.tsv}.

\begin{table}[htbp]
\centering
\small
\setlength{\tabcolsep}{4pt}
\begin{tabular}{lrrrrr}
\hline
Library          & Mean ZVF & Mean last-10 acc & Collapse & Drift & Plateau \\
\hline
GRPO (vanilla)   & 0.481    & 0.379  & 0/5 & 0/5 & 5/5 \\
AERO             & 0.220    & 0.399  & 0/5 & 0/5 & 5/5 \\
CPPO             & 0.295    & 0.392  & 0/5 & 0/5 & 5/5 \\
NGRPO            & 0.300    & 0.387  & 0/5 & 0/5 & 5/5 \\
SCAFGRPO         & 0.317    & 0.409  & 0/5 & 0/5 & 5/5 \\
MCGRPO           & 0.146    & 0.322  & 0/5 & 1/5 & 4/5 \\
GIFT             & 0.114    & 0.326  & 0/5 & 1/5 & 4/5 \\
AREAL            & 0.121    & 0.328  & 0/5 & 1/5 & 4/5 \\
ES               & 0.114    & 0.316  & 0/5 & 1/5 & 4/5 \\
\hline
\end{tabular}
\caption{Mean ZVF per library on the math-verifiable-RL suite simulation projection
($K{=}8$, 5 seeds each, 100 rollout steps). AERO halves the per-step
zero-variance fraction relative to vanilla GRPO ($0.22$ vs $0.48$)
and ties SCAFGRPO on last-10 accuracy ($0.399$ vs $0.409$).
Source: \texttt{experiments/results/zvf_by_library.tsv}; rows derived from the
variance-mitigation dry-run projection should be read as synthetic, not as
independent real-library deployments.}
\label{tab:zvf-by-library}
\end{table}

\paragraph{Reading the table.}
Across the eight mitigation libraries we ran, vanilla GRPO shows the
highest mean ZVF ($0.48$, half of rollouts land in zero-contrast groups).
AERO (arXiv 2602.14338) cuts this in half to $0.22$ and is the
\emph{only} library that ties vanilla GRPO's last-10 accuracy while
halving its gradient starvation. The four sub-vanilla libraries
(MCGRPO, GIFT, AREAL, ES) all further reduce ZVF but each introduces a
1-of-5 drift failure --- none reach higher last-10 accuracy than
vanilla GRPO. CPPO, NGRPO, and SCAFGRPO sit in between on both axes.

\paragraph{Honest scope.}
This cross-library comparison shares \emph{rollout workers, reward
parser, KL coefficient, and template} across all nine libraries, so
differences in Table~\ref{tab:zvf-by-library} reflect the mitigation
mechanism, not confounding infrastructure. It does \emph{not} cover
Dr. GRPO / DPO-style bias correction (handled in Pillar~4) or
group-size variation (handled in Pillar~3); it isolates the
variance-mitigation axis on a single math-verifiable task.

\subsection{What ZVF Predicts Above and Beyond Mean Reward}
\label{sec:zvf-residualised}
The naive correlation between mean ZVF and last-10 accuracy
($r{=}0.22$ across pooled rows; bootstrap 95\% CI $[-0.29,0.94]$,
$n{=}23$) is \emph{insignificant} once the tool-use collapse endpoints
are pooled with arithmetic runs --- mean reward already explains most
of the axis. The cleaner question is whether ZVF carries
\emph{residual} signal after regressing out mean reward.

The residual diagnostic that exists on disk gives the opposite caution:
after regressing out mean reward, residualised ZVF remains \emph{positively}
associated with last-10 accuracy within the variance-mitigation slice
($r=+0.80$, 95\% CI $[+0.54,+1.00]$;
\texttt{experiments/results/zvf\_iter26\_residual.tsv}). We therefore do
not treat ZVF as an incremental predictor above reward mean; the full pooled
correlation breakdown is in \texttt{experiments/results/zvf\_failure\_correlation.tsv}.

\paragraph{Limit.} Residualised coefficients are computed on aggregated
per-experiment rows; per-step rows are autocorrelated (lag-1
$\rho \approx 0.9$) and were deliberately aggregated first.

\paragraph{The clearest case of ZVF-beyond-reward: the all-correct wall.}
A two-seed matched-budget panel (2{,}560 rollouts per arm:
$G{=}2\times160$ steps vs.\ $G{=}16\times20$ steps, Qwen3-8B/GSM8K,
LoRA rank~4, batch~8) supplies the sharpest instance of ZVF carrying
information mean reward cannot. Late in the $G{=}2$ arms, mean train
reward sits near 1.0 --- by the reward axis alone, training is
succeeding --- while $\mathrm{ZVF}$ sits at $0.75$--$1.0$: nearly every
group is all-correct and contributes zero gradient. The $G{=}16$ arms at
the same rollout budget hold $\mathrm{ZVF}\le0.25$. Reward reports
\emph{where the policy is}; ZVF reports \emph{whether training is still
moving} --- and at the high-accuracy end the two disagree completely.
Artifacts: \texttt{tinker-runs/results/er2b\_*.json}.

\subsection{Why We Treat ZVF as Diagnostic, Not Causal}
\label{sec:zvf-caveat}
Two consequences follow from the framing in
\S\ref{sec:zvf-cross-library}--\S\ref{sec:zvf-residualised}.
\begin{itemize}
\item \textbf{Symptom, not cause.} ZVF rises \emph{after} the reward
plateau in every collapse we measured (Nemotron-120B, tool-use
$0\%$, MCGRPO seed 0); it is not the proximal cause of collapse, but
the earliest cheap diagnostic we have for it. On this corpus the first-passage
lead is short: ZVF saturates above $0.8$ before the collapse flag trips, but
only by $1$--$3$ steps in the trajectory-level lead-time artifacts
(\texttt{experiments/results/zvf\_leadtime\_summary.tsv};
\texttt{experiments/results/zvf\_dynamics\_leadtime.tsv}). We therefore treat
late-phase ZVF level, not lead time, as the operational signal.
\item \textbf{Stack-controlled confound.} AERO's halving of ZVF on our
runtime (Table~\ref{tab:zvf-by-library}) is a real, same-stack
measurement; we do \emph{not} claim it transfers to other RL
stacks without re-measurement. The opposite-direction gap between
MCGRPO ($0.15$) and vanilla GRPO ($0.48$) suggests further
variance-mitigation headroom exists.
\end{itemize}

\paragraph{Recommended use.} Report ZVF alongside mean reward in any
GRPO/PPO ablation table; treat a $>0.3$ rise in mean ZVF over
$10$ consecutive rollout steps as a \emph{warning} (not a failure
trigger) that the rollout distribution may be collapsing onto a
single mode for the current prompt batch.

\subsection{Cross-References and Artifacts}
\label{sec:zvf-artifacts}
All numbers and figures cited above are regenerated from
\texttt{scripts/zvf_diagnostic.py}, which itself reads only
\texttt{experiments/results/\{tinker_gsm8k_zvf_*, groupsize_zvf_sweep,
variance_mitigation, tool_code_reward_diagnostics,
scaling_law_three_phase, drgrpo_vs_grpo, samestack_ppo_grpo\}}. The
script's bootstrap-CI math uses $B{=}2000$ paired row resamples; the
deterministic classifier for \texttt{failure_label} is defined at the top
of the file (\S"Failure taxonomy") and re-used unchanged across the
three produced TSVs.

\paragraph{Figures.}
\begin{itemize}
\item \texttt{figures/zvf_vs_failure.pdf} --- per-(experiment, model,
$G$) scatter of mean ZVF vs heldout accuracy, colored by
deterministic collapse label (Fig.~\ref{fig:zvf-vs-failure}, this
section).
\item \texttt{figures/zvf_by_library.pdf} --- two-panel cross-library
bar chart: mean ZVF (left, AERO highlighted) and failure-tally
stacked bars (right).
\end{itemize}

\subsection{Iter 130: Unified ZVF Risk Index}
\label{sec:zvf-risk-index}
We elevate ZVF from a single-axis sanity check to a \emph{multi-channel
risk index} by fusing three orthogonal measurements of the same
underlying rollout distribution:
\begin{enumerate}
\item \textbf{Magnitude channel} --- mean per-step ZVF, mapped through a
logistic centred at $\bar{Z}{=}0.30$ (so the canonical ZVF${>}0.3$ warning
of \S\ref{sec:zvf-cross-library} becomes risk $>0.5$).
\item \textbf{Critical-slowing-down channel} --- rolling lag-1
autocorrelation of the per-step ZVF series (window $w{=}15$, the
Scheffer-2009 EWS statistic), logistic-centred at $0.45$ so that the
GRPO-collapse signature of iter126 (rolling lag-1 $\approx 0.61$ vs
AERO plateau $\approx 0.33$) is the discriminator.
\item \textbf{Drift channel} --- first-half trajectory slope of ZVF,
logistic-centred at $0.004$/step so that GRPO's $9.9{\times}10^{-3}$/step
slope (24$\times$ AERO's $4.1{\times}10^{-4}$/step) registers as risk.
\end{enumerate}

\paragraph{Composite score.} The three channels are combined in two
ways: a \emph{weighted blend} $\mathrm{zvf\_risk} = 0.3\,r_{\rm mag}
+ 0.5\,r_{\rm csd} + 0.2\,r_{\rm drift}$ (CSD-weighted because it is
empirically the strongest single channel), and a \emph{max-fusion}
$\mathrm{zvf\_risk\_max} = \max(r_{\rm mag}, r_{\rm csd}, r_{\rm drift})$
that captures the logical-OR structure: a trajectory is at risk if
\emph{any} channel fires.

\paragraph{Discrimination (cross-experiment).} Across the 52-row
cross-experiment panel (45 variance-mitigation seeds $+$ 2 tool-use
collapse anchors $+$ 5 scaling-law collapse anchors; 11 failures,
41 safe), the per-axis Mann-Whitney AUROCs are:
\begin{center}
\small
\begin{tabular}{lcc}
\hline
Axis / Index & Cross-experiment AUROC & Within-methods AUROC \\
\hline
Magnitude              & 0.663 [0.37, 0.93] & 0.073 [0.00, 0.26] \\
CSD (rolling lag-1)    & \textbf{0.758} [0.52, 0.94] & 0.335 [0.11, 0.57] \\
Drift slope            & 0.579 [0.32, 0.83] & 0.293 [0.09, 0.54] \\
zvf\_risk (weighted)   & 0.763 [0.57, 0.94] & 0.396 [0.21, 0.59] \\
zvf\_risk\_max (OR)    & \textbf{0.929} [0.83, 1.00] & \textbf{0.805} [0.62, 0.95] \\
\hline
\end{tabular}
\end{center}
The max-fusion dominates every single axis (95\% CI excludes 0.5 by a
wide margin on both panels). It is the first ZVF-based diagnostic to
clear the conventional ``substantial'' AUROC threshold ($0.8$) on both
the within-method and the cross-experiment panels.

\paragraph{Why a logical-OR, not a weighted blend?} The within-method
AUROC of magnitude is $0.073$ (anti-discriminating: low-ZVF methods
like GIFT/AREAL/ES fail by collapsing to zero contrast, while GRPO
fails by saturating contrast), and the drift-slope within-method AUROC
is $0.293$ (drift slope does not anticipate GIFT/AREAL-style
crashes). A weighted blend that averages these inverted signals across
seeds washes the discrimination out. The max-fusion is robust because
\emph{either} high ZVF magnitude \emph{or} high CSD \emph{or} high
drift slope suffices to flag the trajectory: this matches the
qualitative observation that GRPO-style collapse, tool-use perfect-zero
collapse, and scaling-law collapse-phase collapse each surface in
\emph{different} channels but never in none of them.

\paragraph{Method-level ranking (max-fusion).} Sorting methods by mean
$\mathrm{zvf\_risk\_max}$ yields the expected collapse-safe ordering:
\begin{center}
\small
\begin{tabular}{lrr}
\hline
Method & $\mathrm{zvf\_risk\_max}$ mean & Failure rate \\
\hline
scaling\_law\_Qwen3-8B        & 0.69 (anchor) & 1.0 \\
scaling\_law\_Nemotron-120B   & 0.69 (anchor) & 1.0 \\
scaling\_law\_Llama-3.1-8B    & 0.69 (anchor) & 1.0 \\
scaling\_law\_Qwen3.5-4B      & 0.69 (anchor) & 1.0 \\
scaling\_law\_DeepSeek-V3.1   & 0.69 (anchor) & 1.0 \\
tool\_use\_llama-8b-inst      & 0.55 (anchor) & 1.0 \\
tool\_use\_qwen3-32b          & 0.55 (anchor) & 1.0 \\
GRPO                          & 0.53          & 0.0 \\
NGRPO / CPPO / AERO           & 0.43--0.49    & 0.0 \\
SCAFGRPO                      & 0.21          & 0.0 \\
GIFT / AREAL / ES             & 0.13--0.20    & 0.2 \\
\hline
\end{tabular}
\end{center}
Note that GRPO is correctly ranked \emph{above} every other
variance-mitigation method on ZVF risk --- even though it does not
collapse on the math-verifiable task --- because its high-CSD,
high-drift profile makes it the closest cousin of the
collapse-positive anchors. AERO, by contrast, is correctly ranked
mid-pack: it has low magnitude and low drift (matching its
collapse-suppression claim from iter118) but moderate CSD ($0.39$
rolling lag-1).

\paragraph{Practitioner use.} Report the max-fusion ZVF Risk Index
alongside mean reward. A trajectory with
$\mathrm{zvf\_risk\_max} > 0.55$ is in the failure-positive regime of
both the within-method and cross-experiment calibration; below
$0.30$ it is in the safe regime. Treat the gap $[0.30, 0.55]$ as
``investigate'': enough of one channel has fired to warrant stopping
the rollout to inspect the per-prompt reward distribution. This is the
first iteration of ZVF as a \emph{three-channel, real-time} diagnostic
rather than a static post-hoc summary.

\paragraph{Artifacts (iter 130).}
\begin{itemize}
\item \texttt{scripts/zvf\_diagnostic\_iter130.py} --- reproducer.
\item \texttt{experiments/results/zvf\_iter130\_risk\_index.tsv} ---
52 rows: \texttt{(method, seed, mean\_zvf, lag1\_zvf\_rolling\_w15,
slope, risk\_mag, risk\_csd, risk\_drift, zvf\_risk, zvf\_risk\_max,
failure\_bin)}.
\item \texttt{experiments/results/zvf\_iter130\_axis\_aurocs.tsv} ---
per-axis + per-scope AUROCs with bootstrap-CI (B=2000).
\item \texttt{experiments/results/zvf\_iter130\_method\_risk.tsv} ---
method-aggregated risk profile.
\item \texttt{experiments/results/zvf\_iter130\_meta.json} ---
machine-readable summary with full per-method breakdown.
\item \texttt{figures/zvf\_vs\_failure.pdf} --- 4-panel
regenerated figure: (a) magnitude $\times$ CSD scatter with failure
colour and anchor markers; (b) per-method risk bar; (c) ROC curves
for all five indices; (d) risk-bin calibration.
\end{itemize}

% Berkeley harvest row 22 (HippoRAG, SP25 L3 Yu Su) -- retrieval-scaffold
% stabilizer; evidence: experiments/results/berkeley/hipporag_eval.tsv
A HippoRAG-style KG+PPR retrieval scaffold~\citep{gutierrez2024hipporag}
over this artifact corpus (67 passages, 21-entity knowledge graph)
matches BM25's recall with a 60\% smaller working set
(recall@10 $= 0.667 =$ BM25 recall@25) and beats random retrieval
decisively (MRR $0.554$ vs.\ $0.248$), but does not beat BM25
head-to-head at this corpus scale (MRR $0.554$ vs.\ $0.569$). At
$N \approx 67$ artifacts, flat lexical retrieval is sufficient; we
retain the scaffold as documentation infrastructure rather than adopt
it as a default
(\texttt{experiments/results/berkeley/hipporag\_eval.tsv}).
